% Source: https://artofproblemsolving.com/community/c3h2918242_argentina_nmo_round_2_intercollegiate_geo__level_1_200616
% Stable identity: aops-topic-2918242-q07
$CA=CB=CD$, so triangle $ACD$ is isosceles with $\angle CDA=\angle CAD=40^\circ$ and $\angle ACD=100^\circ$. Hence $\angle ACB=100^\circ-60^\circ=40^\circ$, and the base angles of triangle $ABC$ are $70^\circ$.
\begin{asy}
usepackage("amsmath");
size(3cm);
defaultpen(fontsize(10pt));
pair C=(0,0),B=dir(-20),A=dir(-60);pair D=rotate(60)*B;draw(A--B--D--C--cycle);draw(C--B,dashed);draw(A--D,dotted);
label("$A$",A,S);label("$B$",B,E);label("$C$",C,W);label("$D$",D,N);
\end{asy}

The angles of $ABDC$ are $\boxed{\angle A=70^\circ,\ \angle B=70^\circ+60^\circ=130^\circ,\ \angle D=60^\circ,\ \angle C=100^\circ}$ (sum $360^\circ$).
